
\documentclass{article}
\usepackage{colortbl}
\usepackage{makecell}
\usepackage{multirow}
\usepackage{supertabular}

\begin{document}

\newcounter{utterance}

\centering \large Interaction Transcript for game `cladder', experiment `full\_v1.5\_default', episode 1169 with Qwen3.5{-}4B.
\vspace{24pt}

{ \footnotesize  \setcounter{utterance}{1}
\setlength{\tabcolsep}{0pt}
\begin{supertabular}{c@{$\;$}|p{.15\linewidth}@{}p{.15\linewidth}p{.15\linewidth}p{.15\linewidth}p{.15\linewidth}p{.15\linewidth}}
    \# & \multicolumn{2}{c}{Player} && \multicolumn{2}{c}{Game Master} \\
    \hline

    \theutterance \stepcounter{utterance}  
    & & & \multicolumn{4}{p{0.6\linewidth}}{
        \cellcolor[rgb]{0.9,0.9,0.9}{
            \makecell[{{p{\linewidth}}}]{
                \texttt{\tiny{[P1$\langle$GM]}}
                \texttt{You are an expert in causal inference. The following question is not a typical commonsense query, but rather a meticulously designed question created by a professor specializing in causal inference, intended to assess the students' mastery of the course content.} \\
\texttt{Imagine a self{-}contained, hypothetical world with only the following conditions, and without any unmentioned factors or causal relationships: Pre{-}conditions has a direct effect on vaccination and disease. Vaccination has a direct effect on disease. The overall probability of getting the vaccine is 59\%. The probability of vaccine refusal and recovering from the disease is 24\%. The probability of getting the vaccine and recovering from the disease is 20\%. Is the chance of recovering from the disease smaller when observing getting the vaccine?} \\
\texttt{Start your answer with 'Yes' or 'No', followed by additional reasoning or evidence to support your explanation.} \\
            }
        }
    }
     \\ \\

    \theutterance \stepcounter{utterance}  
    & & \multicolumn{4}{p{0.6\linewidth}}{
        \cellcolor[rgb]{0.95,0.95,0.95}{
            \makecell[{{p{\linewidth}}}]{
                \texttt{\tiny{[P1$\leftrightarrow$P]}}
                \texttt{{[}DRAFT{]}} \\
\texttt{Yes} \\
\\ 
\texttt{**Reasoning:**} \\
\\ 
\texttt{To determine whether the chance of recovering from the disease is smaller when observing vaccination, we must compare the **marginal probability** of recovery (the overall chance of recovery in the population) with the **conditional probability** of recovery given vaccination (\$P(\textbackslash{}text\{Recovery\} \textbackslash{}mid \textbackslash{}text\{Vaccination\})\$).} \\
\\ 
\texttt{Let's define the events:} \\
\texttt{*   \$V\$: Getting the vaccine} \\
\texttt{*   \$R\$: Recovering from the disease} \\
\texttt{*   \$\textbackslash{}neg V\$: Not getting the vaccine (Vaccine refusal)} \\
\texttt{*   \$\textbackslash{}neg R\$: Not recovering (getting the disease)} \\
\\ 
\texttt{**Step 1: Calculate the Marginal Probability of Recovery (\$P(R)\$)**} \\
\texttt{We are given the probabilities for specific joint scenarios:} \\
\texttt{1.  \$P(\textbackslash{}neg V \textbackslash{}cap R) = 24\textbackslash{}\% = 0.24\$ (Refusal and Recovery)} \\
\texttt{2.  \$P(V \textbackslash{}cap R) = 20\textbackslash{}\% = 0.20\$ (Vaccination and Recovery)} \\
\\ 
\texttt{Since these are the only two mutually exclusive ways to recover in this hypothetical world (either you refused the vaccine and recovered, or you got it and recovered), the total probability of recovery is the sum of these two:} \\
\texttt{\$\$P(R) = P(\textbackslash{}neg V \textbackslash{}cap R) + P(V \textbackslash{}cap R)\$\$} \\
\texttt{\$\$P(R) = 0.24 + 0.20 = 0.44\$\$} \\
\texttt{So, the overall chance of recovering is **44\%**.} \\
\\ 
\texttt{**Step 2: Calculate the Conditional Probability of Recovery Given Vaccination (\$P(R \textbackslash{}mid V)\$)**} \\
\texttt{We are given:} \\
\texttt{*   \$P(V) = 59\textbackslash{}\% = 0.59\$} \\
\texttt{*   \$P(V \textbackslash{}cap R) = 20\textbackslash{}\% = 0.20\$} \\
\\ 
\texttt{Using the definition of conditional probability:} \\
\texttt{\$\$P(R \textbackslash{}mid V) = \textbackslash{}frac\{P(V \textbackslash{}cap R)\}\{P(V)\}\$\$} \\
\texttt{\$\$P(R \textbackslash{}mid V) = \textbackslash{}frac\{0.20\}\{0.59\} \textbackslash{}approx 0.3389\$\$} \\
\texttt{So, the chance of recovering given that one got the vaccine is approximately **33.9\%**.} \\
\\ 
\texttt{**Step 3: Compare the Probabilities**} \\
\texttt{*   Overall chance of recovery (\$P(R)\$): **44\%**} \\
\texttt{*   Chance of recovery given vaccination (\$P(R \textbackslash{}mid V)\$): **\textasciitilde{}33.9\%**} \\
\\ 
\texttt{Since \$33.9\textbackslash{}\% < 44\textbackslash{}\%\$, observing that someone got the vaccine is associated with a lower probability of recovery compared to the general population average.} \\
\\ 
\texttt{**Causal Interpretation (Confounding):**} \\
\texttt{In causal inference terms, this result indicates the presence of **confounding**. The variable "Pre} \\
            }
        }
    }
    & \\ \\

    \theutterance \stepcounter{utterance}  
    & & \multicolumn{4}{p{0.6\linewidth}}{
        \cellcolor[rgb]{0.95,0.95,0.95}{
            \makecell[{{p{\linewidth}}}]{
                \texttt{\tiny{[P1$\leftrightarrow$P]}}
                \texttt{{[}SELF{-}FEEDBACK{]}} \\
\texttt{FORMAT: ok} \\
\texttt{RULES: ok} \\
\texttt{STRATEGY: ok} \\
\texttt{REASON: The answer correctly identifies the confounding bias by comparing the marginal recovery rate (44\%) with the conditional rate given vaccination (33.9\%), proving the chance is indeed smaller for the vaccinated group.} \\
            }
        }
    }
    & \\ \\

    \theutterance \stepcounter{utterance}  
    & \multicolumn{4}{p{0.6\linewidth}}{
        \cellcolor[rgb]{0.561,0.737,0.561}{
            \makecell[{{p{\linewidth}}}]{
                \texttt{\tiny{[P1$\rangle$GM]}}
                \texttt{Yes} \\
\\ 
\texttt{**Reasoning:**} \\
\\ 
\texttt{To determine whether the chance of recovering from the disease is smaller when observing vaccination, we must compare the **marginal probability** of recovery (the overall chance of recovery in the population) with the **conditional probability** of recovery given vaccination (\$P(\textbackslash{}text\{Recovery\} \textbackslash{}mid \textbackslash{}text\{Vaccination\})\$).} \\
\\ 
\texttt{Let's define the events:} \\
\texttt{*   \$V\$: Getting the vaccine} \\
\texttt{*   \$R\$: Recovering from the disease} \\
\texttt{*   \$\textbackslash{}neg V\$: Not getting the vaccine (Vaccine refusal)} \\
\texttt{*   \$\textbackslash{}neg R\$: Not recovering (getting the disease)} \\
\\ 
\texttt{**Step 1: Calculate the Marginal Probability of Recovery (\$P(R)\$)**} \\
\texttt{We are given the probabilities for specific joint scenarios:} \\
\texttt{1.  \$P(\textbackslash{}neg V \textbackslash{}cap R) = 24\textbackslash{}\% = 0.24\$ (Refusal and Recovery)} \\
\texttt{2.  \$P(V \textbackslash{}cap R) = 20\textbackslash{}\% = 0.20\$ (Vaccination and Recovery)} \\
\\ 
\texttt{Since these are the only two mutually exclusive ways to recover in this hypothetical world (either you refused the vaccine and recovered, or you got it and recovered), the total probability of recovery is the sum of these two:} \\
\texttt{\$\$P(R) = P(\textbackslash{}neg V \textbackslash{}cap R) + P(V \textbackslash{}cap R)\$\$} \\
\texttt{\$\$P(R) = 0.24 + 0.20 = 0.44\$\$} \\
\texttt{So, the overall chance of recovering is **44\%**.} \\
\\ 
\texttt{**Step 2: Calculate the Conditional Probability of Recovery Given Vaccination (\$P(R \textbackslash{}mid V)\$)**} \\
\texttt{We are given:} \\
\texttt{*   \$P(V) = 59\textbackslash{}\% = 0.59\$} \\
\texttt{*   \$P(V \textbackslash{}cap R) = 20\textbackslash{}\% = 0.20\$} \\
\\ 
\texttt{Using the definition of conditional probability:} \\
\texttt{\$\$P(R \textbackslash{}mid V) = \textbackslash{}frac\{P(V \textbackslash{}cap R)\}\{P(V)\}\$\$} \\
\texttt{\$\$P(R \textbackslash{}mid V) = \textbackslash{}frac\{0.20\}\{0.59\} \textbackslash{}approx 0.3389\$\$} \\
\texttt{So, the chance of recovering given that one got the vaccine is approximately **33.9\%**.} \\
\\ 
\texttt{**Step 3: Compare the Probabilities**} \\
\texttt{*   Overall chance of recovery (\$P(R)\$): **44\%**} \\
\texttt{*   Chance of recovery given vaccination (\$P(R \textbackslash{}mid V)\$): **\textasciitilde{}33.9\%**} \\
\\ 
\texttt{Since \$33.9\textbackslash{}\% < 44\textbackslash{}\%\$, observing that someone got the vaccine is associated with a lower probability of recovery compared to the general population average.} \\
\\ 
\texttt{**Causal Interpretation (Confounding):**} \\
\texttt{In causal inference terms, this result indicates the presence of **confounding**. The variable "Pre} \\
            }
        }
    }
    & & \\ \\

    \theutterance \stepcounter{utterance}  
    & & & \multicolumn{2}{p{0.3\linewidth}}{
        \cellcolor[rgb]{0.9,0.9,0.9}{
            \makecell[{{p{\linewidth}}}]{
                \texttt{\tiny{[GM$|$GM]}}
                \texttt{Yes} \\
            }
        }
    }
    & & \\ \\

    \theutterance \stepcounter{utterance}  
    & & & \multicolumn{2}{p{0.3\linewidth}}{
        \cellcolor[rgb]{0.9,0.9,0.9}{
            \makecell[{{p{\linewidth}}}]{
                \texttt{\tiny{[GM$|$GM]}}
                \texttt{yes} \\
            }
        }
    }
    & & \\ \\

    \theutterance \stepcounter{utterance}  
    & & & \multicolumn{2}{p{0.3\linewidth}}{
        \cellcolor[rgb]{0.9,0.9,0.9}{
            \makecell[{{p{\linewidth}}}]{
                \texttt{\tiny{[GM$|$GM]}}
                \texttt{game\_result = WIN} \\
            }
        }
    }
    & & \\ \\

\end{supertabular}
}

\end{document}
